\documentclass[twocolumn]{aastex631}
\begin{document}
\title{A residual reionization ionizing-photon-budget shortfall for star-forming galaxies at z\~{}8 (highC systematic corner)}
\author{NebulaMind Lab (autonomous pipeline)}
\affiliation{NebulaMind Lab, \url{https://lab.nebulamind.net}}
\begin{abstract}
Reionization ionizing-photon-budget reconciliation at z\~{}8: star-forming galaxies require f\_esc=0.362 (+0.339/-0.178) to close the budget (Madau-Dickinson SFRD, log xi\_ion=25.5+/-0.15, clumping C=2-5, JWST-SFRD tail), versus indirect-proxy-inferred f\_esc=0.062 (+0.108/-0.039) from LzLCS O32/beta calibrations. Median delta(required-inferred)=+0.274 dex-frac (16-84\%: +0.078 to +0.615); 92\% of the systematic MC shows a shortfall. Conclusion: a genuine SHORTFALL remains, and the sign holds under both O32 and beta calibrations. The result is bounded by the xi\_ion x clumping x proxy-calibration systematic, not by statistics. Generated autonomously from public data (jwst) via the NebulaMind Lab runner: a bounded, reproducible, descriptive study.
\end{abstract}
\section{Introduction}
Recent studies have highlighted a potential crisis in understanding the photon budget required for reionization [Muñoz2024]. This has led to increased demands on ionizing sources, suggesting that current models may not fully account for the necessary photons to drive reionization [Davies2021]. To address this issue, we revisit the ionizing-photon-budget calculation using a literature-anchored approach. Our work builds upon previous efforts to calibrate excursion set reionization models and assess the galaxy ionizing photon budget at high redshifts [Park2022, Duncan2015].
\section{Data and method}
We adopt the Madau \& Dickinson (2014) analytic fitting function for the cosmic star formation rate density (SFRD) and utilize published values for xi\_ion and O32/beta f\_esc proxy calibrations from LzLCS studies [Chisholm+22, Flury+22; Simmonds+24]. Our method focuses on reconciling the reionization ionizing-photon-budget at z\~{}8 using these literature-anchored values.
\section{Result}
Our calculation reveals that star-forming galaxies require an escape fraction of f\_esc=0.362 (+0.339/-0.178) to close the budget, considering the Madau-Dickinson SFRD, log xi\_ion=25.5±0.15, clumping factor C=2-5, and JWST-SFRD tail. However, indirect-proxy-inferred f\_esc from LzLCS O32/beta calibrations yields a significantly lower value of 0.062 (+0.108/-0.039). The median delta between required and inferred values is +0.274 dex-frac (16-84\%: +0.078 to +0.615), with 92\% of systematic Monte Carlo simulations showing a shortfall.
\begin{figure}\centering\includegraphics[width=\columnwidth]{result.png}\caption{A residual reionization ionizing-photon-budget shortfall for star-forming galaxies at z\~{}8 (highC systematic corner)}\end{figure}
\section{Caveats}
It is essential to acknowledge the limitations of our approach. Our analysis relies on automated, single-selection, uncalibrated measurements, which may introduce biases and uncertainties. The result is bounded by systematics in xi\_ion, clumping factor, and proxy-calibration, rather than statistical errors. Further research is needed to refine these parameters and improve our understanding of the reionization process. Additionally, incorporating data from future surveys and observations will be crucial for validating and refining our findings.
\section*{References}
\footnotesize
[Muoz2024] Muñoz et al. (2024). Reionization after JWST: a photon budget crisis?. bibcode 2024MNRAS.535L..37M \par [Davies2021] Davies et al. (2021). The Predicament of Absorption-dominated Reionization: Increased Demands on Ionizing Sources. bibcode 2021ApJ...918L..35D \par [Park2022] Park et al. (2022). Calibrating excursion set reionization models to approximately conserve ionizing photons. bibcode 2022MNRAS.517..192P \par [Duncan2015] Duncan et al. (2015). Powering reionization: assessing the galaxy ionizing photon budget at z < 10. bibcode 2015MNRAS.451.2030D \par [Madau2017] Madau (2017). Cosmic Reionization after Planck and before JWST: An Analytic Approach. bibcode 2017ApJ...851...50M

\end{document}
